\input{../tex-inputs/latex-leseansicht-vorspann}

\section[Gustav Schwarzkopf an Arthur Schnitzler, 26. 6. 1900]{L04299 Gustav Schwarzkopf an Arthur Schnitzler, 26. 6. 1900}
\nopagebreak\mylabel{L04299v}
\rehead{ }\normalsize\beginnumbering\briefempfaengerindex{Schnitzler, Arthur@\textsc{Schnitzler, Arthur}!zzzSchwarzkopf, Gustav@\emph{von Gustav Schwarzkopf}!1900-06-262@{26. 6. 1900}|(be}
\toendnotes[C]{\smallbreak\pagebreak[2]}
\correspDesc{Versand  durch Gustav Schwarzkopf am 26. 6. 1900 in Hinterbrühl
\newline{}Übermittlung  am 26. 6. 1900 in Hinterbrühl
\newline{}Zustellung  am 27. 6. 1900 in Wien
\newline{}Erhalt  durch Arthur Schnitzler am 27. 6. 1900 in Wien}\toendnotes[C]{\smallbreak}
\Standort{CUL, Schnitzler, B 96a.}
\physDesc{Postkarte, 332 Zeichen
\newline{}Handschrift: schwarze Tinte, deutsche Kurrent
\newline{}Versand: 1) Stempel: »\nobreak{}\oindex{Hinterbrühl@\textbf{Hinterbrühl}, \emph{Hauptstadt}|pwk}Hinterbrühl, 26. 6. 00, 6–7N\nobreak{}«.   2) Stempel: »\nobreak{}\oindex{IX., Alsergrund@\textbf{IX., Alsergrund}, \emph{Verwaltungsgebiet}|pwk}Wien 9/3 72, 27. 6. 00, 1.N, Bestellt\nobreak{}«. }\toendnotes[C]{\smallbreak}\pstart{}{\pb}Herrn \textsc{D\textsuperscript{r} Arthur Schnitzler}\pend{}\pstart{}\textsc{Wien\oindex{Wien@\textbf{Wien}, \emph{Verwaltungsgebiet}|pw}}\pend{}\pstart{}IX. Frankgaſſe 1\oindex{Wien@\textbf{Wien}!IX., Alsergrund@\textbf{IX., Alsergrund}!Frankgasse 1@\textbf{Frankgasse 1}, \emph{Wohngebäude}|pw}.\pend{}{\bigskip}\vspace{1em}
\pstart
           \raggedleft{}{\pb}26. VI 900\pend
           \vspace{0.5em}
\pstart
           Lieber Arthur, man redet mir{ }ſehr zu noch einige Tage hier zu
               bleiben und da Sie \label{K_L04299-1v}\edtext{morgen nach Baden\oindex{Baden bei Wien@\textbf{Baden bei Wien}, \emph{Hauptstadt}|pw}}{\lemma{\textnormal{\emph{morgen nach Baden}}}\Cendnote{\textnormal{Vgl. A. S.: \emph{Tagebuch}, 27. 6. 1900.}}}\label{K_L04299-1} gehen,
               würden wir uns ja doch kaum{ }ſehen. Vor der ganz großen Ferienreiſe ko{\geminationm}en Sie doch wol noch einmal von Reichenau\oindex{Reichenau an der Rax@\textbf{Reichenau an der Rax}, \emph{Verwaltungsgebiet}|pw} nach Wien\oindex{Wien@\textbf{Wien}, \emph{Verwaltungsgebiet}|pw}. Nicht?
               Bis dahin gute Unterhaltung.\pend
           
\pstart
           Herzlichſt{\\[\baselineskip]}Ihr{\\[\baselineskip]}\spacefill\mbox{Gustav.}\pend
           \leftskip=0em{}\selectlanguage{ngerman}\endnumbering\briefempfaengerindex{Schnitzler, Arthur@\textsc{Schnitzler, Arthur}!zzzSchwarzkopf, Gustav@\emph{von Gustav Schwarzkopf}!1900-06-262@{26. 6. 1900}|)be}\mylabel{L04299h}
\begin{anhang}
\end{anhang}\newcommand{\dateiname}{L04299}\newcommand{\titel}{Gustav Schwarzkopf an Arthur Schnitzler, 26. 6. 1900}\newcommand{\editorInnen}{Herausgegeben von Jahnke, SelmaMüller, Martin Anton}\input{../tex-inputs/latex-leseansicht-abspann}
